% Nano-drones, MCU, ASIC, neuromorphique
\section{Gamme nano : MCU, ASIC et neuromorphique}
\label{sec:nano}

\subsection{Le GAP8 et la lignée PULP}

Le \glslink{gap8}{GAP8} \citep{2018-flamand-gap-8-risc-v-soc} est la puce de l'\glslink{aideck}{AI-deck} du \glslink{crazyflie}{Crazyflie},
utilisée par la plupart des travaux nano du corpus. C'est un \gl{SoC}{SoC} RISC-V de 55\,nm : un
contrôleur et un cluster de 8 cœurs, 512\,kB de L2, 64\,kB de L1, un accélérateur de
convolution, environ 10\,\glslink{mac}{GMAC/s} en crête et au plus 75\,\glslink{u-mw}{mW} au pire cas \src{§II,
Table~II}. Le papier, court et écrit par le concepteur, ne rapporte aucune détection :
seulement CIFAR-10 en 650\,µs pour 15\,mW, méthode de mesure non décrite \src{§IV}. Son
enseignement durable est que la gestion explicite de la mémoire (\gl{tuilage}{tuilage} L3/L2/L1 par
\gl{DMA}{DMA}) est le cœur de l'effort de portage \deduit, comme sur un accélérateur \gl{FPGA}{FPGA} \gl{dataflow}{dataflow}.

Le tableau~\ref{tab:nano} rassemble les réseaux déployés sur \glslink{gap8}{GAP8}, sur son successeur \glslink{gap9}{GAP9}
(section~\ref{sec:gap9}) et sur \gl{MCU}{MCU}.

\begin{table}[!htbp]
\centering\footnotesize
\caption{Réseaux déployés en gamme nano. Tâches très différentes : les chiffres donnent
des ordres de grandeur. Les deux lignes GAP9 n'ont aucune puissance mesurée, et la
latence de Tiny-DroNeRF est simulée.}
\label{tab:nano}
\begin{tabularx}{\textwidth}{@{}P{3.0cm}LP{2.0cm}rP{2.4cm}@{}}
\toprule
Travail & Tâche, réseau & Puce & Débit & Puissance (\glslink{portee}{portée}), source \\
\midrule
\glslink{pulp-dronet}{PULP-Dronet} \citep{arxiv-1805.01831-64mw-dnn-visual-navigation-engine}\pp & navigation, \glslink{dronet}{DroNet}
  Fixed16, 41\,\glslink{mac}{MMAC} & \glslink{gap8}{GAP8} & 6 / 18\,\glslink{fps}{FPS} & 64 / 284\,\glslink{u-mw}{mW} (shield), §VI\med \\
\glslink{tiny-pulp-dronet}{Tiny-PULP-Dronet} \citep{arxiv-2407.02405-tiny-pulp-dronets-squeezing-neural}\pp & navigation,
  1,5\,\glslink{mac}{MMAC}, 6,4\,kB & \glslink{gap8}{GAP8} & 89 / 160\,\glslink{fps}{FPS} & 34\,\glslink{u-mw}{mW} (chip), 0,38\,\glslink{u-mj}{mJ}/image, §III-B\med \\
\glslink{pulp-frontnet}{PULP-Frontnet} \citep{arxiv-2103.10873-fully-onboard-ai-powered-human}\pp & pose humaine, \gl{INT8}{INT8},
  4\,\glslink{mac}{MMAC} & \glslink{gap8}{GAP8} & 134,7\,\glslink{fps}{FPS} & 86,6\,\glslink{u-mw}{mW} (chip), §V-B\med \\
\glslink{nanoflownet}{NanoFlowNet} \citep{arxiv-2209.06918-nanoflownet-real-time-dense-optical}\pp & \gl{flot}{flot optique} dense,
  INT8 & \glslink{gap8}{GAP8} & 5,57\,\glslink{fps}{FPS} & \nr, Table~IV\med \\
\citet{arxiv-2301.12175-bio-inspired-autonomous-exploration-policies}\pp & détection \glslink{ssd-mobilenetv2}{SSD-MbV2} 1,0$\times$,
  2 classes, 534\,\glslink{mac}{MMAC} ; 50\,\% \glslink{map}{mAP} & \glslink{gap8}{GAP8} & 1,6\,\glslink{fps}{FPS} & 134,5\,\glslink{u-mw}{mW} (\glslink{aideck}{AI-deck}), Table~II \\
\glslink{blankskip}{BlankSkip} \citep{arxiv-2603.28149-blankskip-early-exit-object-detection}\pp & même détecteur +
  \gl{sortie}{sortie anticipée} ; 0,540 \glslink{map}{mAP} & \glslink{gap8}{GAP8} & 1,86\,\glslink{fps}{FPS} moyen & \nr, Table~2 \\
\citet{arxiv-2603.02936-self-supervised-domain-adaptation-visual}\pp & pose 3D d'une porte de course,
  CNN INT8, 318\,\glslink{mac}{MMAC} ; \glslink{mae}{MAE} $\psi$ 13,1\textdegree & \glslink{gap9}{GAP9 (NE16)} & 32,8\,\glslink{fps}{FPS} & \nr, §V, Table~I\med \\
\glslink{tiny-dronerf}{Tiny-DroNeRF} \citep{arxiv-2603.01850-tiny-dronerf-tiny-neural-radiance}\pp & \emph{entraînement}
  \glslink{nerf}{NeRF}, FP16, 160$\times$160 ; \glslink{psnr}{PSNR} 26,94\,dB & \glslink{gap9}{GAP9} (cluster) & 73\,\glslink{u-ms}{ms}/pas, simulé & 60\,\glslink{u-mw}{mW} (non décrite), §III-B, §V-C\med \\
\glslink{tinyissimoyolo}{TinyissimoYOLO} \citep{arxiv-2306.00001-tinyissimoyolo-quantized-low-memory-footprint} & détection
  88$\times$88, 3 classes \gl{VOC}{VOC} ; 58,5\,\% \glslink{map}{mAP} & \glslink{max78000}{MAX78000} & 5,5\,\glslink{u-ms}{ms} & 196\,\glslink{u-uj}{µJ}/inférence (MCU), §III\med \\
\glslink{tinyissimoyolo}{TinyissimoYOLO} & idem & \glslink{stm32h7a3}{STM32H7A3} & 359\,\glslink{u-ms}{ms} & \nr, §III\med \\
\bottomrule
\end{tabularx}
\end{table}

\paragraph{Navigation.} \glslink{pulp-dronet}{PULP-Dronet} \citep{arxiv-1805.01831-64mw-dnn-visual-navigation-engine}\pp{}
(publié dans IEEE IoT Journal d'après le DOI) est le premier \gl{CNN}{CNN} de navigation en boucle
fermée entièrement à bord d'un nano-drone de 27\,g. À 64\,\glslink{u-mw}{mW}, le calcul représente 0,8\,\%
de la puissance du drone, estimée à 7,6\,\glslink{u-w}{W} \src{Fig.~12}. Son apport le plus utile ici est
système : l'autonomie en vol stationnaire passe de 440\,\glslink{u-s}{s} pour le \glslink{crazyflie}{Crazyflie} seul à 340\,s
avec le shield qui calcule, et la masse du module en coûte l'essentiel (90\,s contre 10\,s
pour le calcul) \src{Table~IV}. \glslink{tiny-pulp-dronet}{Tiny-PULP-Dronet}
\citep{arxiv-2407.02405-tiny-pulp-dronets-squeezing-neural}\pp{} réduit le réseau d'un
facteur 50 en taille et 27 en \glslink{mac}{MAC}, pour un gain de latence de 8,5 seulement : la première
couche et le format mémoire deviennent le goulot \src{§III-B}. Aucun vol n'est rapporté,
et le multitâche qui motive le travail n'est pas démontré.

\paragraph{Pose et flot optique.} \glslink{pulp-frontnet}{PULP-Frontnet}
\citep{arxiv-2103.10873-fully-onboard-ai-powered-human}\pp{} ferme la boucle de contrôle
sur une estimation de pose humaine. Son résultat le plus instructif : le réseau le plus
rapide (135\,\glslink{u-hz}{Hz}) suit moins bien la personne que le plus lent (48\,Hz), avec une erreur de
position médiane de 0,63 contre 0,41\,m \src{Table~IV}. Le débit maximal n'est pas la
bonne cible \deduit. \glslink{nanoflownet}{NanoFlowNet}
\citep{arxiv-2209.06918-nanoflownet-real-time-dense-optical}\pp{} calcule un flot optique
dense à 5,57\,\glslink{fps}{FPS} sur \glslink{gap8}{GAP8} et sert l'évitement d'obstacles, sans mesure de puissance ni
précision après \gl{quantif}{quantification} \src{Table~IV}.

\paragraph{Détection.} Deux articles seulement font de la détection sur \glslink{gap8}{GAP8}, sur la même
tâche jouet : deux classes d'objets posés au sol dans une pièce.
\citet{arxiv-2301.12175-bio-inspired-autonomous-exploration-policies}\pp{} obtiennent
1,6\,\glslink{fps}{FPS} pour 134,5\,\glslink{u-mw}{mW} à l'\glslink{aideck}{AI-deck} \src{Table~II} et un \glslink{detrate}{taux de détection} de 90\,\% en
mission à 0,5\,m/s, qui tombe à 83\,\% à 1\,m/s \src{Table~III}. Le bilan de puissance du
drone entier est un chiffre clé : moteurs 7,32\,\glslink{u-w}{W} (91\,\%), \glslink{aideck}{AI-deck} 0,134\,W (1,67\,\%)
\src{Table~IV}. \glslink{blankskip}{BlankSkip}
\citep{arxiv-2603.28149-blankskip-early-exit-object-detection}\pp{} ajoute une branche qui
détecte les images vides et arrête l'inférence : +24\,\% de débit moyen pour $-$0,015 de
\glslink{map}{mAP} \src{Table~2}. Le gain dépend de la proportion d'images vides, fabriquée par les
auteurs sur le jeu de test, et la latence au pire cas reste celle du modèle complet
\deduit.

\paragraph{MCU à accélérateur.} \glslink{tinyissimoyolo}{TinyissimoYOLO}
\citep{arxiv-2306.00001-tinyissimoyolo-quantized-low-memory-footprint} montre que
l'accélérateur CNN intégré au MCU fait la différence : 5,5\,\glslink{u-ms}{ms} sur \glslink{max78000}{MAX78000} contre 359\,ms
sur \glslink{stm32h7a3}{STM32H7A3}, pour le même réseau \src{§III}. Sa \glslink{map}{mAP} (3 classes, au plus 3 objets par
image, entrée 88$\times$88) ne se compare à aucune autre du corpus, et varie de 57 à
59\,\% selon la section du papier.

\subsection{La génération GAP9}
\label{sec:gap9}

Trois preprints récents remplacent l'\glslink{aideck}{AI-deck} par le \glslink{gap9shield}{GAP9Shield}, module enfichable sur le
\glslink{crazyflie}{Crazyflie} construit autour du \glslink{gap9}{GAP9}, successeur du \glslink{gap8}{GAP8} chez le même fondeur. Le papier de
présentation \citep{arxiv-2407.13706-gap9shield-150gops-ai-capable-ultra}\pp{} décrit un
cluster RISC-V de 9 cœurs avec l'accélérateur \glslink{gap9}{NE16} et un calcul flottant FP32/FP16/bfloat16,
jusqu'à 370\,\glslink{u-mhz}{MHz}, 1,6\,MB de L2 (contre 512\,kB pour le \glslink{gap8}{GAP8}) et 330\,µW/GOP
\src{§2.1}. Ce sont des chiffres constructeur, et ils ne s'accordent pas d'un papier à
l'autre : 150\,\glslink{gops}{GOPS} dans le titre, mais 15,6\,\glslink{gops}{GOPS} (\glslink{dsp}{DSP}) et 32,2\,\glslink{mac}{GMAC/s} (ML) dans le texte
\src{§2.1} ; 150\,GOp/s entiers pour le seul \glslink{gap9}{NE16} selon
\citet[§III-B]{arxiv-2603.01850-tiny-dronerf-tiny-neural-radiance}\pp{} ; 220\,GOp/s
combinés pour moins de 100\,\glslink{u-mw}{mW} selon
\citet[§IV-A]{arxiv-2603.02936-self-supervised-domain-adaptation-visual}\pp. Le shield
ajoute une caméra OV5647 de 5\,MP, un module Wi-Fi et cinq capteurs de distance ToF, pour
environ 6\,g, contre 7\,g pour l'\glslink{aideck}{AI-deck} et le Multi-ranger réunis \src{§3.1}. Le texte
annonce une réduction de masse de 15\,\% (§3.1) puis de 20\,\% (Abstract) ; le rapport
6/7 donne 14\,\% \deduit. Ses seules mesures propres portent sur les périphériques, sans
méthode décrite : caméra 80\,mW en QVGA à 15\,\glslink{fps}{FPS} et 120\,mW en VGA à 45\,\glslink{fps}{FPS}, Wi-Fi
250--260\,mW, capteurs ToF au plus 150\,mW \src{§3.1}. La promesse « sous 100\,mW » ne vaut
donc que pour le SoC en calcul, pas pour le module \deduit. Les chiffres de détection
qu'il met en avant (\glslink{yolo}{YOLO} en 17\,\glslink{u-ms}{ms} et 1,59\,\glslink{u-mj}{mJ}) viennent de travaux cités, sans précision
ni résolution, et ne sont pas mesurés dans le papier \src{§3.2} ; ils n'entrent pas dans
le tableau~\ref{tab:nano}. Schémas et PCB sont annoncés, pas publiés à la lecture
\src{Intro}.

\paragraph{Inférence sur le NE16.}
\citet{arxiv-2603.02936-self-supervised-domain-adaptation-visual}\pp{} estiment la pose
3D d'une porte de course (sous-thème vol agile) avec un CNN entraîné en simulation, puis
adapté au réel sans étiquette : une perte de cohérence entre les poses prédites et
l'odométrie embarquée du drone remplace la capture de mouvement. Quantifié en INT8 et
déployé sur le \glslink{gap9}{NE16}, le réseau de 318\,\glslink{mac}{MMAC} tourne en 30,4\,\glslink{u-ms}{ms}, ou 32,8\,\glslink{fps}{FPS} \src{§V},
soit environ 10,4\,\glslink{mac}{GMAC/s} effectifs, contre 0,85 pour le \glslink{ssd-mobilenetv2}{SSD-MobileNetV2} sur \glslink{gap8}{GAP8}
(534\,\glslink{mac}{MMAC} à 1,6\,\glslink{fps}{FPS}) : un ordre de grandeur, entre réseaux et chaînes d'outils
différents \deduit.
L'erreur sur le jeu réel reste de 26, 28 et 10\,cm en position et 13\textdegree{} en lacet
\src{Table~I}\med, pour une porte de 100$\times$80\,cm. Trois limites : la puissance du
modèle déployé n'est pas mesurée (\nr ; les 100\,\glslink{u-mw}{mW} sont ceux du constructeur), la
précision après quantification est dite « sans perte » sans être chiffrée, et aucun vol en
boucle fermée n'est rapporté \src{§VI}.

\paragraph{Apprentissage à bord.} \glslink{tiny-dronerf}{Tiny-DroNeRF}
\citep{arxiv-2603.01850-tiny-dronerf-tiny-neural-radiance}\pp{} (accepté à ICRA 2026
selon les auteurs) est le seul travail du corpus qui \emph{entraîne} un réseau sur un
nano-drone : un \glslink{nerf}{NeRF} dérivé d'\glslink{instantngp}{Instant-NGP}, réduit à 21,4\,MB contre 527\,MB, entraîné en
FP16 sur le cluster par rétropropagation \src{§I, §IV-B}. Le \glslink{gap9}{NE16} reste inutilisé, car il
ne sait faire que la passe avant en entier \src{§III-B}. Un pas d'entraînement prend
73\,\glslink{u-ms}{ms}, soit 97\,min pour une reconstruction, et le rendu d'une image 160$\times$160 environ
4,1\,\glslink{u-s}{s} \src{§V-C}, avec un \glslink{psnr}{PSNR} de 26,94\,dB, 5,7\,dB sous \glslink{instantngp}{Instant-NGP}, mais avec des
images de 160 contre 800\,px de côté \src{Table~I}\med. Ces latences sortent du simulateur
GVSOC, pas de la puce \src{§V-C}. La puissance de 60\,\glslink{u-mw}{mW} est une « moyenne » à 370\,\glslink{u-mhz}{MHz}
sans méthode décrite \src{§III-B}, et ni l'énergie d'un entraînement de 97\,min ni son
effet sur une batterie de 350\,mAh ne sont chiffrés. En conditions réelles, le \glslink{psnr}{PSNR} tombe
à 21,2\,dB, et l'entraînement de cette expérience n'est pas dit fait à bord \src{§V-D}.
L'apprentissage fédéré sur un essaim de quatre drones est simulé, sans canal radio réel
\src{§V-E}.

Le \glslink{gap9}{GAP9} élargit donc la \glslink{gamme}{gamme nano} vers le flottant et l'entraînement, et multiplie la
mémoire sur puce par trois, sans changer d'ordre de grandeur de puissance d'après le
constructeur. Mais aucun des trois papiers ne mesure la puissance d'un réseau sur \glslink{gap9}{GAP9}, ni
à la puce ni à la carte, et aucun ne fait de détection : ce que la génération \glslink{gap9}{GAP9} apporte
à un détecteur embarqué reste à mesurer \deduit.

\subsection{ASIC et neuromorphique}

\paragraph{Navion.} \citet{2019-suleiman-navion-2-mw-fully-integrated} intègrent une
odométrie visuelle-inertielle complète sur un \gl{ASIC}{ASIC} de 65\,nm, sans DRAM externe : 24\,\glslink{u-mw}{mW} en
configuration maximale et 2\,mW à 20\,\glslink{fps}{FPS} quand les paramètres sont adaptés à chaque
séquence (puissance de la puce) \src{§V-A, §V-D}\med. La méthode de co-design est
transposable à un détecteur : tout sur puce, tampons d'image compressés, parcimonie
exploitée \deduit. Les 2\,mW supposent un réglage fait a posteriori, et aucun vol n'est
rapporté : les données sont rejouées depuis un \glslink{zc706}{ZC706} \src{§V-B}.

\paragraph{Loihi.} Deux travaux de contrôle de drone utilisent la puce neuromorphique \glslink{loihi}{Loihi}
d'Intel. \citet{arxiv-2108.03694-event-driven-vision-control-uavs}\pp{} calculent une
transformée de Hough sur les événements d'une caméra \gl{DVS}{DVS} en 250\,µs, soit une estimation à
20\,kHz contre 1\,kHz sur \gl{CPU}{CPU}, sur un banc à un degré de liberté \src{Fig.~2, §III-A}\med ;
aucune puissance n'est mesurée. \citet{arxiv-2303.08778-fully-neuromorphic-vision-control-autonomous}\pp{}
ferment la boucle vision $\rightarrow$ contrôle sur un quadrirotor de 1\,kg, avec un \gl{SNN}{SNN} de
flot optique à 27\,\glslink{u-uj}{µJ} par inférence à 200\,\glslink{u-hz}{Hz} \src{Abstract}. Deux réserves : l'énergie
est mesurée sans entrées-sorties, alors que l'USB est le vrai goulot selon les auteurs, et
la comparaison à trois ou quatre ordres de grandeur oppose \glslink{loihi}{Loihi} à une émulation du SNN en
PyTorch sur \glslink{jetson-nano}{Jetson Nano}, ce que les auteurs reconnaissent \src{Discussion}. Aucun des deux
ne fait de détection d'objets.

\subsection{Ce que la gamme nano apprend au reste}

La \glslink{gamme}{gamme nano} est la seule du corpus où la part du calcul dans le budget du drone est
mesurée : 1,67\,\% de la puissance pour la vision
\citep[Table~IV]{arxiv-2301.12175-bio-inspired-autonomous-exploration-policies}, et un
coût en autonomie dominé par la masse
\citep[Table~IV]{arxiv-1805.01831-64mw-dnn-visual-navigation-engine}. Elle montre aussi que
le débit limite la vitesse de mission (\glslink{detrate}{taux de détection} qui baisse à 1\,m/s) et que la
qualité de contrôle ne croît pas avec le débit (\glslink{pulp-frontnet}{PULP-Frontnet}). Ces trois constats n'ont
pas d'équivalent mesuré pour un FPGA ou une \glslink{jetson}{Jetson} sur un drone de taille moyenne. Ils
reposent tous sur le \glslink{gap8}{GAP8} : pour le \glslink{gap9}{GAP9}, le corpus n'a encore ni puissance mesurée ni
bilan à l'échelle du drone (section~\ref{sec:gap9}).
